% 图契约：本题的源、边界和推导逐项对应；箭头符号按图注解释。
\documentclass[tikz,border=7pt]{standalone}
\usepackage{fontspec,xeCJK,amsmath,amssymb}
\IfFontExistsTF{Times New Roman}{\setmainfont{Times New Roman}}{\setmainfont{TeX Gyre Termes}}
\IfFontExistsTF{Heiti SC}{\setCJKmainfont[AutoFakeBold=1.5]{Heiti SC}}{\setCJKmainfont[AutoFakeBold=1.5]{FandolHei-Regular}}
\usetikzlibrary{arrows.meta,bending,calc,patterns,positioning}
\definecolor{ink}{HTML}{243745}
\definecolor{blue}{HTML}{275DAB}
\definecolor{green}{HTML}{177C63}
\definecolor{red}{HTML}{BD4935}
\definecolor{orange}{HTML}{AA6C16}
\definecolor{purple}{HTML}{715393}
\tikzset{every node/.style={font=\fontsize{9.3}{12.5}\selectfont,text=ink},
axis/.style={-{Stealth[length=4pt]},line width=.65pt,draw=ink!65},
vec/.style={-{Stealth[length=5pt]},line width=1pt,draw=blue},
short/.style={-{Stealth[length=3pt]},line width=.8pt},
guide/.style={dashed,draw=ink!40,line width=.5pt},
card/.style={rounded corners=2pt,draw=blue!65,fill=blue!4,align=center,text width=6.5cm,inner xsep=9pt,inner ysep=7pt},
title/.style={font=\fontsize{10.5}{14}\selectfont\bfseries,text width=15cm,align=center},
label/.style={font=\fontsize{8.4}{11}\selectfont,align=center},
note/.style={label,text width=6.2cm}}
\begin{document}
\begin{tikzpicture}[x=1cm,y=1cm]
\path[use as bounding box] (-8,-3.5) rectangle(8,3.5);
\node[title] at (0,3.15) {M3-015\quad 稀疏介质球混合物的等效电容率};
\begin{scope}[shift={(-4,0)}]
\draw[guide] (-2.1,-1.5) rectangle(2.1,1.5);
\foreach \x/\y in {-1.4/-.8,-1.0/.8,0/0,1.3/.9,1.4/-.85}{
\filldraw[draw=ink!50,fill=blue!7] (\x,\y)circle(.32);
\draw[vec,green] (\x-.22,\y)--(\x+.22,\y);}
\node[label] at(0,1.8){宏观体积 $\Delta V$，颗粒数 $n\Delta V$};
\draw[vec] (-2,-2)--(2,-2) node[midway,below] {$\mathbf E$};
\node[label] at(0,-1.2){颗粒半径 $a$；$f\ll1$};\end{scope}
\node[card,draw=blue!65,fill=blue!4] at (3.9,1.8) {$\displaystyle\mathbf p=4\pi\varepsilon_0a^3\frac{\varepsilon_r-1}{\varepsilon_r+2}\mathbf E$};
\node[card,draw=green!65,fill=green!4] at (3.9,0) {$\mathbf P=n\mathbf p+O(f^2\varepsilon_0E)$};
\node[card,draw=orange!65,fill=orange!4] at (3.9,-1.8) {$\displaystyle\varepsilon_{r,\rm eff}=1+3f\frac{\varepsilon_r-1}{\varepsilon_r+2}+O(f^2)$};
\node[label,text width=15cm] at(0,-3.05){箭头示 $\varepsilon_r>1$；若 $0<\varepsilon_r<1$，各颗粒的诱导偶极矩与外场反向。};
\end{tikzpicture}\end{document}